% GENERATED FILE. Edit canonical lessons, then run npm run generate:books.
% canonical-source-hash: fnv1a64:ce20eec738685a79
% canonical-lessons: TE-C172-yavvanam, TE-C172-vrddhapyam, TE-C172-intalo, TE-C172-eta, TE-C172-eppatiki

\chapter{Youth, Old age, Meanwhile, Every year, Forever}
\label{ch:te-youth-old-age-meanwhile-every-year-forever}
\hlchaptermodality{172}

\begin{chapteropening}
\textbf{By the end of this chapter:} \emph{I can say youth, old age, meanwhile, every year and forever.}

Complete the last lesson of chapter 172: I can say youth, old age, meanwhile, every year and forever.
\end{chapteropening}

\section[\texorpdfstring{\te{యవ్వనం} (yavvanaṁ)}{యవ్వనం (yavvanaṁ)}]{\te{యవ్వనం} (yavvanaṁ) — youth}

\label{lesson:TE-C172-yavvanam}



\begin{quote}
Before the new one: say the Telugu for rarely, then the Telugu for constantly.
\end{quote}

\subsection*{You'll want to know: \te{యవ్వనం}}
\textbf{\te{యవ్వనం}} — \emph{yavvanaṁ} — \textquotedblleft{}youth\textquotedblright{}.

\begin{grammarlens}[title={What you've built}]
One more word for this chapter.
\end{grammarlens}

\subsection*{Your turn}
\begin{itemize}
\raggedright
  \item \emph{Say it:} \emph{yavvanaṁ}
  \item \emph{Say it:} \emph{yavvanaṁ}, once more
  \item \emph{Say it:} say \emph{yavvanaṁ}
  \item \emph{Recall:} say the Telugu for rarely, then the Telugu for constantly, then say \emph{yavvanaṁ} again
\end{itemize}

\begin{tcolorbox}[breakable,colback=teal!4,colframe=teal!35!black,title={Before you move on}]
What does \te{యవ్వనం} mean? (\textquotedblleft{}Youth\textquotedblright{}.) Say it once more.
\end{tcolorbox}
\section[\texorpdfstring{\te{వృద్ధాప్యం} (vṛddhāpyaṁ)}{వృద్ధాప్యం (vṛddhāpyaṁ)}]{\te{వృద్ధాప్యం} (vṛddhāpyaṁ) — old age}

\label{lesson:TE-C172-vrddhapyam}



\begin{quote}
Before the new one: say the Telugu for constantly, then the Telugu for youth.
\end{quote}

\subsection*{You'll want to know: \te{వృద్ధాప్యం}}
\textbf{\te{వృద్ధాప్యం}} — \emph{vṛddhāpyaṁ} — \textquotedblleft{}old age\textquotedblright{}.

\begin{grammarlens}[title={What you've built}]
One more word for this chapter.
\end{grammarlens}

\subsection*{Your turn}
\begin{itemize}
\raggedright
  \item \emph{Say it:} \emph{vṛddhāpyaṁ}
  \item \emph{Say it:} \emph{vṛddhāpyaṁ}, once more
  \item \emph{Say it:} say \emph{vṛddhāpyaṁ}
  \item \emph{Recall:} say the Telugu for constantly, then the Telugu for youth, then say \emph{vṛddhāpyaṁ} again
\end{itemize}

\begin{tcolorbox}[breakable,colback=teal!4,colframe=teal!35!black,title={Before you move on}]
What does \te{వృద్ధాప్యం} mean? (\textquotedblleft{}Old age\textquotedblright{}.) Say it once more.
\end{tcolorbox}
\section[\texorpdfstring{\te{ఇంతలో} (intalō)}{ఇంతలో (intalō)}]{\te{ఇంతలో} (intalō) — meanwhile}

\label{lesson:TE-C172-intalo}



\begin{quote}
Before the new one: say the Telugu for youth, then the Telugu for old age.
\end{quote}

\subsection*{You'll want to know: \te{ఇంతలో}}
\textbf{\te{ఇంతలో}} — \emph{intalō} — \textquotedblleft{}meanwhile\textquotedblright{}.

\begin{grammarlens}[title={What you've built}]
One more word for this chapter.
\end{grammarlens}

\subsection*{Your turn}
\begin{itemize}
\raggedright
  \item \emph{Say it:} \emph{intalō}
  \item \emph{Say it:} \emph{intalō}, once more
  \item \emph{Say it:} say \emph{intalō}
  \item \emph{Recall:} say the Telugu for youth, then the Telugu for old age, then say \emph{intalō} again
\end{itemize}

\begin{tcolorbox}[breakable,colback=teal!4,colframe=teal!35!black,title={Before you move on}]
What does \te{ఇంతలో} mean? (\textquotedblleft{}Meanwhile\textquotedblright{}.) Say it once more.
\end{tcolorbox}
\section[\texorpdfstring{\te{ఏటా} (ēṭā)}{ఏటా (ēṭā)}]{\te{ఏటా} (ēṭā) — every year, annually}

\label{lesson:TE-C172-eta}



\begin{quote}
Before the new one: say the Telugu for old age, then the Telugu for meanwhile.
\end{quote}

\subsection*{You'll want to know: \te{ఏటా}}
\textbf{\te{ఏటా}} — \emph{ēṭā} — \textquotedblleft{}every year, annually\textquotedblright{}.

\begin{grammarlens}[title={What you've built}]
One more word for this chapter.
\end{grammarlens}

\subsection*{Your turn}
\begin{itemize}
\raggedright
  \item \emph{Say it:} \emph{ēṭā}
  \item \emph{Say it:} \emph{ēṭā}, once more
  \item \emph{Say it:} say \emph{ēṭā}
  \item \emph{Recall:} say the Telugu for old age, then the Telugu for meanwhile, then say \emph{ēṭā} again
\end{itemize}

\begin{tcolorbox}[breakable,colback=teal!4,colframe=teal!35!black,title={Before you move on}]
What does \te{ఏటా} mean? (\textquotedblleft{}Every year, annually\textquotedblright{}.) Say it once more.
\end{tcolorbox}
\section[\texorpdfstring{\te{ఎప్పటికీ} (eppaṭikī)}{ఎప్పటికీ (eppaṭikī)}]{\te{ఎప్పటికీ} (eppaṭikī) — forever}

\label{lesson:TE-C172-eppatiki}



\begin{quote}
Before the new one: say the Telugu for meanwhile, then the Telugu for every year.
\end{quote}

\subsection*{You'll want to know: \te{ఎప్పటికీ}}
\textbf{\te{ఎప్పటికీ}} — \emph{eppaṭikī} — \textquotedblleft{}forever\textquotedblright{}.

\begin{grammarlens}[title={What you've built}]
One more word for this chapter.
\end{grammarlens}

\subsection*{Your turn}
\begin{itemize}
\raggedright
  \item \emph{Say it:} \emph{eppaṭikī}
  \item \emph{Say it:} \emph{eppaṭikī}, once more
  \item \emph{Say it:} say \emph{eppaṭikī}
  \item \emph{Recall:} say the Telugu for meanwhile, then the Telugu for every year, then say \emph{eppaṭikī} again
\end{itemize}

\begin{tcolorbox}[breakable,colback=teal!4,colframe=teal!35!black,title={Before you move on}]
What does \te{ఎప్పటికీ} mean? (\textquotedblleft{}Forever\textquotedblright{}.) Say it once more.
\end{tcolorbox}
